Projekt został zrealizowany na podstawie modelu \emph{klient-serwer} \figref{fig:clientserver}, polegającego na rozproszeniu aplikacji na dwa główne komponenty:
\begin{itemize}
    \item \textbf{backend} (serwer) -- odpowiedzialny za przetwarzanie żądań klientów, obsługę logiki biznesowej oraz komunikację z zewnętrznymi usługami,
    \item \textbf{frontend} (aplikacja kliencka) -- interaktywny interfejs, który wysyła zapytania do serwera oraz prezentuje odpowiednio sformatowane dane użytkownikowi.
\end{itemize}

\begin{figure}[h]
    \centering
    \includegraphics[width=0.8\textwidth]{client-server_model}
    \figcaption{Przykładowy schemat modelu klient-serwer}{\cite{wikipedia_clientserver}}
    \label{fig:clientserver}
\end{figure}

Komponenty komunikują się poprzez sieć za pomocą protokołu HTTP\footnote{Hypertext Transfer Protocol -- protokół komunikacyjny wykorzystywany przez przeglądarki}.
Dostęp do zasobów serwera odbywa się poprzez interfejs programowania aplikacji zgodny z architekturą REST\footnote{Representational State Transfer -- styl architektoniczny identyfikujący elementy systemu jako zasoby}.
Komunikacja odbywa się w sposób \emph{bezstanowy}, czyli serwer nie zapisuje kontekstu komunikacji po swojej stronie.
Wszystkie informacje dotyczące sesji muszą zostać przekazane łącznie z samym zapytaniem przez aplikację kliencką.
Zasoby są identyfikowane jako unikalne adresy URL\footnote{Uniform Resource Locator -- adres zasobu dostępnego w sieci komputerowej}, a operacje na nich są zdefiniowane jako podstawowe metody protokołu HTTP: GET, POST, PATCH, DELETE \autocite{rest_tutorial}.

\section{Serwer}

\subsection{Warstwowy system}

\subsection{Obsługa żądań}

\subsection{Walidacja danych}

\subsection{Autentykacja za pośrednictwem ciasteczek}

\subsection{Rozsyłanie maili}

\subsection{Zarządzanie zdjęciami}

\subsection{Przeliczanie reputacji sprzedawcy}

\subsection{Obsługa błędów}

\subsection{Baza danych NoSQL}

\subsection{Opis użytych modeli danych}
